These results use the independently selected authoritative Hayward cohort. Residuals below are unweighted final residuals; standard tomography fits absolute picks only, so its catalog-difference RMS is an additional diagnostic. Catalog displacement measures relocation, not error against known source positions.

The selected pair graph has 35 connected components; the largest contains 1058 events, and 34 events have no differential links. Unlinked events remain in the inversion with absolute observations. 28576 catalog differences span metadata epochs at identical NSLC and receiver coordinates; this geometry match does not establish unchanged timing calibration.

\begin{center}\begin{tabular}{lrrrrr}\toprule Method & $R_{\rm abs}$ (s) & $R_{\rm cat}$ (s) & Updates & LSQR & Min. rays (\%)\\\midrule
Initial & 0.39163 & 0.09340 & -- & -- & --\\
standard & 0.07869 & 0.07162 & 12/13 & 13/13 & 95.75\\
dd & 0.16064 & 0.06160 & 18/18 & 18/18 & 95.76\\
loc1d & 0.29188 & 0.25423 & 10/13 & 13/13 & 97.93\\
\bottomrule\end{tabular}\end{center}
The initial row uses the first pre-update forward prediction. Updates counts accepted steps; LSQR counts solves meeting the existing relative normal-residual tolerance of $10^{-8}$. The last column is the minimum converged fraction over all used rays in any saved pre-update history row. These history entries precede their update and do not represent the final model. Frozen-weight objectives are compared within each update; weights can change between iterations. The solver stops at its iteration limit or when no line-search step is accepted; it does not test nonlinear stationarity or a small-source-update criterion.

No waveform-correlation observations are included. The history column for precise-difference RMS is therefore empty (NaN), not a nonfinite modeled residual.

For standard, median and maximum catalog displacements are 1.100 and 5.956 km; 3 sources lie within $10^{-6}$ km of a grid boundary. Nodal velocities range from 3.326 to 6.288 km/s. There are 1 zero-update iterations and 0 accepted objective increases above the existing $10^{-8}$ audit tolerance; the worst LSQR relative normal residual is 5.14e-09. The final recorded iteration changes a source coordinate by at most 0.000 km and nodal log slowness by at most 0. Geoid depths span -0.960 to 16.494 km, with 1 negative-depth estimates and 0 estimates outside the grid beyond the boundary tolerance. Boundary counts by face are 1 at $y=-30$ km; 2 at $y=30$ km; a source on intersecting faces contributes to each count. Across all 14262--14262 rays per recorded iteration, the largest projected/raw gradient norms are 0.00577/0.00577 s/km. At the final model, production/refined convergence counts are 31/31 of 32 sampled rays, with maximum projected gradients 0.00049/0.000399 s/km and raw gradients 0.00049/0.000399 s/km. Halving the segment spacing changes sampled travel times by at most 0.101 ms. For 24 source-coordinate derivatives, maximum production/refined errors against centered travel-time differences are 2.88e-05/2.16e-05 s/km; 0/0 components exceed the existing 0.003 s/km component tolerance. 0 derivative components are omitted at domain boundaries.

For dd, median and maximum catalog displacements are 1.178 and 8.799 km; 15 sources lie within $10^{-6}$ km of a grid boundary. Nodal velocities range from 4.427 to 5.618 km/s. There are 0 zero-update iterations and 0 accepted objective increases above the existing $10^{-8}$ audit tolerance; the worst LSQR relative normal residual is 6.97e-09. The final recorded iteration changes a source coordinate by at most 2.761 km and nodal log slowness by at most 0.00541. Geoid depths span -1.000 to 17.244 km, with 19 negative-depth estimates and 0 estimates outside the grid beyond the boundary tolerance. Boundary counts by face are 2 at $y=30$ km; 13 at $z=-1$ km; a source on intersecting faces contributes to each count. Across all 14262--14262 rays per recorded iteration, the largest projected/raw gradient norms are 0.00544/0.00544 s/km. At the final model, production/refined convergence counts are 32/32 of 32 sampled rays, with maximum projected gradients 8.08e-07/2.77e-07 s/km and raw gradients 8.08e-07/2.77e-07 s/km. Halving the segment spacing changes sampled travel times by at most 0.008 ms. For 24 source-coordinate derivatives, maximum production/refined errors against centered travel-time differences are 1.52e-06/1.24e-07 s/km; 0/0 components exceed the existing 0.003 s/km component tolerance. 0 derivative components are omitted at domain boundaries.

For loc1d, median and maximum catalog displacements are 2.638 and 8.350 km; 8 sources lie within $10^{-6}$ km of a grid boundary. Nodal velocities range from 3.770 to 7.180 km/s. There are 3 zero-update iterations and 0 accepted objective increases above the existing $10^{-8}$ audit tolerance; the worst LSQR relative normal residual is 9.4e-09. The final recorded iteration changes a source coordinate by at most 0.000 km and nodal log slowness by at most 0. Geoid depths span -1.000 to 15.282 km, with 23 negative-depth estimates and 0 estimates outside the grid beyond the boundary tolerance. Boundary counts by face are 1 at $y=30$ km; 7 at $z=-1$ km; a source on intersecting faces contributes to each count. Across all 14262--14262 rays per recorded iteration, the largest projected/raw gradient norms are 0.00795/0.00795 s/km. At the final model, production/refined convergence counts are 32/32 of 32 sampled rays, with maximum projected gradients 9.32e-06/9.84e-06 s/km and raw gradients 9.32e-06/9.84e-06 s/km. Halving the segment spacing changes sampled travel times by at most 0.164 ms. For 24 source-coordinate derivatives, maximum production/refined errors against centered travel-time differences are 0.000121/8.83e-05 s/km; 0/0 components exceed the existing 0.003 s/km component tolerance. 0 derivative components are omitted at domain boundaries.

Some source estimates reach the computational top, added to contain elevated receivers. These constrained estimates do not establish source-depth recovery. Negative geoid depth means elevation above the geoid, not necessarily above local terrain. The velocity sections show the complete computational depth range so these estimates remain visible.

Over the saved model nodes, DD minus standard velocity has mean -0.0213 km/s, RMS 0.4164 km/s, and maximum absolute difference 1.3754 km/s. The median separation between their source estimates is 0.623 km. These differences do not establish which field model is more accurate.

The absolute-data resolution calculation is linearized at the final DD model with fixed uncertainty weights, first-difference log-slowness smoothing, and column-scaled damping. The model Schur-complement condition number is 2.81e+03, its inverse relative residual is 4.01e-15, and the smallest damped source-block eigenvalue is 3.39e-05. Resolution diagonals range from -0.00011 to 0.838; conditional velocity-noise standard deviations range from 0.00407 to 0.033 km/s. For this final DD linearization, 14260/14262 unique absolute-data rays meet the ray stopping criterion; maximum projected/raw gradient norms are 2.95e-05/0.000162 s/km.

This is a conditional absolute-data calculation, not full DD resolution. Reported symmetric pick uncertainties are interpreted as independent one-standard-deviation errors; provider confidence levels are absent. Propagation through the damped estimator excludes correlated station/model errors, regularization bias, and nonlinear or ray errors. It eliminates unconstrained source parameters and does not represent the active domain bounds on some inferred source coordinates. The plotted sections interpolate nodal resolution and nodal noise standard deviations; they do not propagate covariance to interpolated velocity. Small propagated noise therefore does not establish small total velocity error. Ray convergence means a projected gradient below $10^{-7}$ s/km, or a travel-time change below $10^{-11}$ s with gradient below $10^{-5}$ s/km. The source finite-difference half-step is 0.01 km. Small spacing changes alone do not establish derivative convergence or globally earliest arrivals, and the deterministic final samples do not certify every production ray.
